\begin{center}2.\end{center}
\noindent\hspace*{\parindent}Die zur Analysis gehörigen Begriffsbestimmungen, ihre Unterscheidungen und gegliederten Classificirungen, die in den Lehrbüchern aufgestellt zu werden pflegen, lassen zum Theil noch erkennen, dass man bei ihrer Einführung sich auf einem niedern Standpunkte mit beschränkterm Gesichtspunkte befand. Man schuf eine Begriffsbildung, um diesem oder jenem praktischen Bedürfnisse entgegenzukommen, und dachte bei jener an keine weitern Grenzen als dieses Bedürfniss erforderte. Über diese Grenzen hinaus hatte die Begriffsbestimmung keinen klaren Sinn mehr. Kein Wunder also, dass wenn man sich auf Fragen einliess, wo die selben Begriffe ausserhalb des Gebiets, in welchem allein sie ihre Berechtigung hatten, in Anspruch genommen wurden, Widersprüche und Verwirrung die Folge waren. Es gehören dahin z. B. die lange streitigen Fragen über die Logarithmen der negativen Grössen, die immer nur ein Streit de lana caprina bleiben mussten, so lange nicht der Begriff von Logarithm aus Einem Guss auf eine für das ganze Gebiet der Grössen gültige, vollkommen klare Art festgestellt war. Zu den Ursachen, welchen man eine solche im 17. und 18. Jahrhundert so oft vorkommende und auch zum Theil noch in das gegenwärtige hineinreichende Unzulänglichkeit bei den mathematischen Begriffsbestimmungen vorzüglich zuzuschreiben hat, und auf deren
\par\vfill\hfill{\small 51\textsuperscript{*}}\par
\clearpage\SourceLeaf{411}\ArithmeticHead{404}{L}{ANALYSIS. NACHLASS.}
\noindent specielle Erörterung ich hier mich nicht einlassen will, gehört auch der verkehrte Gesichtspunkt, aus welchem man die jetzt sogenannten imaginären Grössen so lange betrachtet hat. Wir sehen dieselben zuerst unter dem Namen von unmöglichen Grössen auftreten bei denjenigen quadratischen Gleichungen, wo die Auflösung die Ausziehung der Quadratwurzel aus einer negativen Grösse erfordert. In so weit enthielt der Ausdruck, dass die Wurzeln der Gleichung unmöglich seien, eigentlich nichts weiter als eine Verneinung der Existenz von Wurzeln. Späterhin machte man die Bemerkung, dass auch bei algebraischen Gleichungen von höhern Graden in solchen Fällen, wo sie entweder gar nicht auflösbar waren oder wenigstens nicht eine ihrem Grade gemässe Anzahl von Wurzeln darboten, die Unauflösbarkeit oder Unvollzähligkeit der Auflösungen immer nur von der Unmöglichkeit, aus einer negativen Grösse eine Quadratwurzel zu erhalten, abhing, und \textsc{d’Alembert} und \textsc{Euler} generalisirten diese Bemerkung (dem Wesen nach, wenn auch nicht in dieser klaren Ausdrucksweise) dahin, dass es nur der Zulassung einer fingirten Grösse \(\surd{-1}\) bedürfe, um jeder entwickelten algebraischen Gleichung die ihrer Ordnungszahl gleiche Menge von Wurzeln zu verschaffen. Ein genügender strenger Beweis dies[es] wichtigsten Lehrsatzes in der Theorie der algebraischen Gleichungen ist zwar erst viel später gelungen; aber schon von jener Zeit an wurde es immer mehr üblich, die unschicklichen Benennungen von unmöglichen Grössen fahren zu lassen, und die jene fingirte \(\surd{-1}\) involvirenden Grössen imaginäre zu nennen, im Gegensatz der reellen, unter denen man die Totalität aller positiven und negativen begriff. Seit der Zeit sind die imaginären Grössen, wie eine besondere Gattung von Grössen, in den analytischen Calcul aufgenommen, und man hat sich wohl dabei gestanden, indem ihre Beihülfe sehr häufig auf eine überraschend leichte Art zu einem Ziele führt, welches ohne sie nur viel mühsamer sich erreichen lassen würde, oder indem ihre Zuziehung mancher mathematischen Lehre, die, wenn sie bloss auf das Gebiet der reellen Grössen beschränkt werden müsste, nur ungelenk und lückenhaft erscheinen würde, die schönste Abrundung gibt.

Bei allem dem sind die imaginären Grössen, so lange ihre Grundlage immer nur in einer Fiction bestand, in der Mathematik nicht sowohl wie eingebürgert, als viel mehr nur wie geduldet betrachtet, und weit davon entfernt geblieben, mit den reellen Grössen auf gleiche Linie gestellt zu werden.
\clearpage\SourceLeaf{412}\ArithmeticHead{405}{R}{CONVERGENZ DER ENTWICKLUNG PERIODISCHER FUNCTIONEN.}
\noindent Zu einer solchen Zurücksetzung ist aber jetzt kein Grund mehr, nachdem die Metaphysik der imaginären Grössen in ihr wahres Licht gesetzt und nachgewiesen ist, dass diese, eben so gut wie die negativen, ihre reale gegenständliche Bedeutung haben\textsuperscript{*}). Die vollständige Erkenntniss der Natur einer analytischen Function muss auch die Einsicht in ihr Verhalten bei den imaginären Werthen des Arguments in sich schliessen, und oft ist sogar letztere unentbehrlich zu einer richtigen Beurtheilung der Gebarung der Function im Gebiete der reellen Argumente. Unerlässlich ist es daher auch, dass die ursprüngliche Festsetzung des Begriffs der Function sich mit gleicher Bündigkeit über das ganze Grössengebiet erstrecke, welches die reellen und die imaginären Grössen unter dem gemeinschaftlichen Namen der complexen Grössen in sich begreift[\textsuperscript{**})].
\ClosingRule
\begin{center}[Anhang.]\end{center}
\begin{center}\begin{minipage}{.45\linewidth}
\hrule\smallskip\centering\small [Ein Blatt in Fa, Kapsel {\scriptsize 46}a.]\par\smallskip\hrule
\end{minipage}\end{center}
\begin{center}\bfseries Maasstab für Convergenz der Reihen.\end{center}
\begin{center}Allgemeines Glied der Reihe \(R_n\), der Summenreihe \(S_n\).\end{center}
\noindent\hspace*{\parindent}Es sei \(\rho_nR_n\) das allgemeine Glied einer neuen Reihe, die weder convergent noch divergent ist. Man wird dann die Convergenzen der Reihen \((R)\) und \(\left(\frac1\rho\right)\) als gleich zu betrachten haben.

Es wird \(\rho_n\) in die Form zu bringen sein
\[
\rho_n=e^{\alpha n}n^\beta(\log n)^\gamma\,{.}\,(\log\log n)^\delta\,{.}\,(\log^3 n)^\varepsilon\ \text{u.s.w.},
\]
\noindent wo \(\alpha,\beta,\gamma,\delta,{.}\,{.}\,{.}\) positive, negative Zahlen oder 0 sein können, aber die erste, die nicht verschwindet, jedenfalls positiv.

Die Bedeutung lässt sich so aussprechen:
\par\vfill\noindent\rule{.25\linewidth}{.4pt}\par\smallskip
{\footnotesize
\noindent\hspace*{\parindent}\textsuperscript{*}) Göttingische gelehrte Anzeigen, 1831, St. 64. [Werke II, S. 169.]\par
\noindent\hspace*{\parindent}[\textsuperscript{**}) In der Handschrift folgen noch zwei Artikel, die hier nicht mit abgedruckt wurden.]\par}
\clearpage\SourceLeaf{413}\ArithmeticHead{406}{L}{ANALYSIS. NACHLASS.}
\noindent\hspace*{\parindent}Die Reihe
\begingroup\small\setlength{\arraycolsep}{3pt}\renewcommand{\arraystretch}{1.85}
\[
\begin{array}{@{}l|c|c|c|c@{}}
\begin{array}{@{}l@{}}
e^{\alpha'n}R_n\\\hline
e^{\alpha n}n^{\beta'}R_n\\\hline
e^{\alpha n}n^\beta(\log n)^{\gamma'}R_n\\\hline
e^{\alpha n}n^\beta(\log n)^\gamma(\log\log n)^{\delta'}[R_n]\\\hline
e^{\alpha n}n^\beta(\log n)^\gamma(\log^2 n)^\delta(\log^3 n)^{\varepsilon'}[R_n]
\end{array}
&\renewcommand{\arraystretch}{1.32}\begin{array}{c}\text{conver-}\\\text{girt}\\\text{oder}\\\text{diver-}\\\text{girt}\\\text{je nach}\\\text{dem}\end{array}
&\begin{array}{c}\alpha'\\\hline\beta'\\\hline\gamma'\\\hline\delta'\\\hline\varepsilon'\end{array}
&\begin{array}{c}\text{kleiner}\\\text{oder}\\\text{grösser}\\\text{[ist]}\\\text{als}\end{array}
&\begin{array}{c}\alpha\\\hline\beta\\\hline\gamma\\\hline\delta\\\hline\varepsilon\end{array}
\end{array}
\]
\endgroup
\noindent Wenn eine der Reihen
\[
e^{\alpha n}R_n,\ e^{\alpha n}n^\beta R_n,\ e^{\alpha n}n^\beta(\log n)^\gamma[R_n]\ \text{u.s.w.}
\]
\noindent weder divergirt noch convergirt, so ist der Ausdruck damit zu Ende. Trifft d[ies] z. B. bei der dritten zu, so ist \(\delta=0\), \(\varepsilon=0\), etc.

Die Summenreihe \(S_n\), unter [der] Voraussetzung, dass alle Glieder von \((R)\) wenn nicht gleich, doch von irgend einer Stelle an gleiche Zeichen haben, wird unendlich, wenn die erste der Grössen
\[
\alpha-1,\ \beta-1,\ \gamma-1,\ \text{u.s.w.},
\]
\noindent die nicht \(=0\) wird, negativ ist oder wenn alle \(=0\) [sind].
\ClosingRule
